\section{Parameter Identification}

\begin{frame}{White Box Approach }
    \begin{block}{Physical Parameters Estimation}
        \textbf{Assumptions}: Full \textbf{knowledge of the equations} of the system . \\
        \textbf{Goal}: Identify the \textbf{six parameters vectore} : $ m_1, m_2, l_1, l_2, b_1, b_2 $.
    \end{block}
    
    \begin{figure}
        \centering
        \includegraphics[width=0.6\textwidth]{stock2.png}
    \end{figure}
\end{frame}

\begin{frame}{Physical Parameters Estimation - Architecture}
    

    \textbf{Training Recipe}:
    \begin{itemize}
        \item Loss: \textbf{MSE} between predicted and true derivatives.
        \item \textbf{Curricular Learning}: progressively increasing the integration step $h$ during training.
        \item Data Transformation: trigonometric transformation of angles \textbf{(sin, cos)}.
        \item \textbf{Batching} and Learning Rate Scheduling.
    \end{itemize}

    \begin{figure}
        \centering
        \includegraphics[width=0.6\textwidth]{stock.png}
    \end{figure}

\end{frame}

\begin{frame}{Physical Parameters Estimation - Results}
    \textbf{Observations}:
    \begin{itemize}
        \item \textbf{Gauge Invariance Property}.
        \item \textbf{Parameter Convergence}: optimal Identification.
        \item \textbf{Correct Evolution} even though little deviations of parameters $\rightarrow$ high deviations after a while.
        \item \textbf{Energy evolution}: Perfectly matching (\textit{by-construction}).
        \item \textbf{Metrics}: R2 score, Max Parameter Discrepancy.
        \item \textbf{Fixing Parameters}: helps the identification of the others.
    \end{itemize}
    \begin{figure}
        \centering
        \includegraphics[width=0.8\textwidth]{stock.png}
    \end{figure}
\end{frame}


% \begin{frame}
%     \frametitle{Architecture for Physical Parameters Estimation}

%     We propose to use an integrator plus a dynamic model with learnable physical parameters:

%     \begin{figure}
%         \centering
%         \includegraphics[width=0.4\textwidth]{images/integrator.png}
%     \end{figure}

%     Training recipe:
%     \begin{itemize}
%         \item loss (MSE, MAE)
%         \item curricular steps (increasing h)
%         \item data transformation (sin-cos)
%     \end{itemize}
    
% \end{frame}

% \begin{frame}
%     \frametitle{Architecture for Physical Parameters Estimation}

%     \begin{figure}
%         \centering
%         \includegraphics[width=0.4\textwidth]{images/dyn_ode.png}
%     \end{figure}

%     Fixing some parameters can help the estimation of the others.

%     \begin{figure}
%         \centering
%         \includegraphics[width=0.4\textwidth]{images/dyn_ode_fixed.png}
%     \end{figure}
    
% \end{frame}

% \begin{frame}
%     \frametitle{Results for Physical Parameters Estimation}

%     Fixing parameters make the difference.

%     Gauge invariance issue can affect the estimation.

%     \begin{figure}
%         \centering
%         \includegraphics[width=0.7\textwidth]{images/phys_params_results.png}
%     \end{figure}
% \end{frame}

% \subsection{Dynamics Matrices Entries}
